% Source: https://artofproblemsolving.com/community/c3h2136021_a_random_problem_i_came_up_with
% Stable identity: aops-topic-2136021-q01
\textbf{The condition.} If $n^2-100=\dfrac{7mq}{2}$ with integers $m,q$, then $2\left(n^2-100\right)=7mq$, so $7\mid2\left(n^2-100\right)$, and since $\gcd(7,2)=1$, $7\mid n^2-100$. Conversely, if $n^2-100=7k$ for an integer $k$, take $m=2k$ and $q=1$. So the condition is exactly
\[7\mid n^2-100=(n-10)(n+10).\]
Since $7$ is prime, this means $n\equiv10\equiv3\pmod 7$ or $n\equiv-10\equiv4\pmod7$.

\textbf{Counting.} The primes less than $40$ are
\[2,3,5,7,11,13,17,19,23,29,31,37,\]
twelve in all. Their remainders modulo $7$ are $2,3,5,0,4,6,3,5,2,1,3,2$, so exactly four of them, namely $3,11,17,31$, satisfy the condition. The probability is $\dfrac{4}{12}=\dfrac{1}{3}$, so $a=1$, $b=3$ and
\[a^2b+ab^2+ab+a+2b=3+9+3+1+6=\boxed{22}.\]
